% Generated by supplement/build_r6_paper_table.py from the
% hash-pinned supplement/CLASSIFICATION-RECORDS.json.
\begin{longtable}{c l r r c r l}
\caption{Exceptional redundancy-six classes.  Each row is one
projective-semilinear class.  The columns give a canonical syndrome
representative, the size and stabilizer order of each constituent
\(\PGL_2(q)\)-orbit, the number \(\phi\) of such orbits fused by
coefficient Frobenius, the shared-root pair count \(E\), and the binary-quintic
factor type \(\tau_5\) in odd characteristic.}
\label{tab:r6-exceptional}\\
\toprule
\(q\) & representative & \(|\mathcal O|\) & \(|G_f|\) & \(\phi\) & \(E\) & \(\tau_5\)\\
\midrule
\endfirsthead
\toprule
\(q\) & representative & \(|\mathcal O|\) & \(|G_f|\) & \(\phi\) & \(E\) & \(\tau_5\)\\
\midrule
\endhead
\bottomrule
\endfoot
7 & \([0:0:0:1:0:1]\) & 168 & 2 & 1 & 7 & \(1^3+2\)\\
7 & \([0:0:1:0:3:1]\) & 336 & 1 & 1 & 19 & \(1^2+1+2\)\\
7 & \([0:0:1:0:3:2]\) & 336 & 1 & 1 & 15 & \(1^2+3\)\\
7 & \([0:1:0:0:0:1]\) & 168 & 2 & 1 & 10 & \(1+2+2\)\\
7 & \([0:1:0:0:2:0]\) & 112 & 3 & 1 & 33 & \(1+1+3\)\\
7 & \([0:1:0:1:1:2]\) & 336 & 1 & 1 & 9 & \(1+4\)\\
7 & \([0:1:0:1:1:5]\) & 336 & 1 & 1 & 9 & \(1+1+3\)\\
7 & \([0:1:0:1:3:0]\) & 336 & 1 & 1 & 21 & \(1+1+1+2\)\\
7 & \([0:1:0:1:3:3]\) & 336 & 1 & 1 & 7 & \(1+1+3\)\\
7 & \([0:1:0:1:3:6]\) & 336 & 1 & 1 & 14 & \(1+1+3\)\\
7 & \([0:1:0:3:0:5]\) & 168 & 2 & 1 & 11 & \(1+4\)\\
7 & \([0:1:0:3:0:6]\) & 168 & 2 & 1 & 20 & \(1+4\)\\
7 & \([0:1:0:3:2:1]\) & 336 & 1 & 1 & 13 & \(1+1+3\)\\
7 & \([1:0:0:0:1:2]\) & 336 & 1 & 1 & 7 & \(1+4\)\\
7 & \([1:0:0:0:1:3]\) & 336 & 1 & 1 & 23 & \(5\)\\
7 & \([1:0:0:3:1:3]\) & 336 & 1 & 1 & 13 & \(1+4\)\\
7 & \([1:0:1:0:5:2]\) & 336 & 1 & 1 & 8 & \(1+4\)\\
7 & \([1:0:1:0:6:2]\) & 336 & 1 & 1 & 8 & \(5\)\\
8 & \([0:0:0:1:0:0]\) & 9 & 56 & 1 & 0 & --\\
8 & \([0:1:1:0:2:1]\) & 504 & 1 & 3 & 20 & --\\
8 & \([0:1:1:0:2:5]\) & 504 & 1 & 3 & 19 & --\\
8 & \([1:0:0:1:3:6]\) & 504 & 1 & 3 & 14 & --\\
8 & \([1:0:1:3:5:2]\) & 168 & 3 & 1 & 54 & --\\
9 & \([0:1:0:0:0:4]\) & 180 & 4 & 2 & 28 & \(1+4\)\\
9 & \([0:1:0:1:4:2]\) & 720 & 1 & 2 & 29 & \(1+4\)\\
11 & \([0:0:0:1:0:0]\) & 132 & 10 & 1 & 60 & \(1^3+1^2\)\\
11 & \([0:1:0:2:0:10]\) & 660 & 2 & 1 & 30 & \(1+4\)\\
13 & \([0:1:0:0:0:2]\) & 546 & 4 & 1 & 60 & \(1+4\)\\
\end{longtable}
\noindent For \(q=8\), integers are binary polynomial-basis codes
with \(\alpha^3=\alpha+1\); for \(q=9\), they are ternary
polynomial-basis codes with \(\alpha^2=-1\).  Thus
\(\sum\phi=(18, 11, 4, 2, 1)\) and
\(\sum\phi|\mathcal O|=(5152, 4713, 1800, 792, 546)\),
in the field order \(7,8,9,11,13\).
